\begin{proposition}[Known with derivation: the opened Earle--Sipe coordinate comparison]\label{A1-Earle-Sipe-comparison}
Let $z,w$ be disjoint local unit-disc coordinates on the two branches to be joined. For $0<r_1,r_2<1$ and $t=r_1r_2e^{i\theta}$, plumbing by $zw=t$ is precisely the analytic boundary welding which removes $|z|<r_1$ and $|w|<r_2$ and identifies
\[
 z=r_1e^{i\varphi},\qquad w=r_2e^{i(\theta-\varphi)}.
\]
Its normal tensor is $t\,\partial_z\otimes\partial_w$. This is the same coordinate-level comparison as the original Earle--Sipe construction, with the full local coordinates retained.
\end{proposition}
\begin{proof}
In \cite[\S2, equations (2.1)--(2.2), pp.~80--82]{EarleSipe1991A567}, the two local identifications are $z_1w_1=t$ and $z_2w_2=t^{2p-3}$. Section~6, equations (6.1)--(6.2), p.~88, constructs the identical fibers by removing discs of radii $|t|^{1/2}$ and $|t|^{(2p-3)/2}$, respectively, and identifying their boundary circles by those same equations. These passages are read in the original published article.

For the displayed general radii, set $u=z/r_1$, $v=w/r_2$. The boundary equation becomes $uv=e^{i\theta}$; extending across the analytic circles gives exactly $w=t/z$ on the annular overlap. Moving the cut within that overlap changes neither quotient, by the identity on the retained charts. Finally
\[
 e^{i\theta}\,\partial_u\otimes\partial_v
 =e^{i\theta}r_1r_2\,\partial_z\otimes\partial_w
 =t\,\partial_z\otimes\partial_w.
\]
Thus this opened original supplies a direct comparison with Proposition~\ref{A1-welding-plumbing}. It does not remove the dependence on the coordinate germs or provide an inverse on unrestricted normal tensors.
\end{proof}

\paragraph{What is and is not checked in the joint Earle--Marden source.}
The distinct Earle--Marden chapter \emph{On holomorphic families of Riemann surfaces}, Contemp. Math.~573 (2012), 67--97, must not be substituted for \emph{Holomorphic plumbing coordinates}, Contemp. Math.~575 (2012), 41--52. The latter is exactly reference [24] in the opened \cite{Matsumoto2019A567}. His printed p.~27 reports its \S13, Theorem~II, distinguishes the asserted embedding over a whole Teichm\"uller stratum from Hubbard--Koch's local product, and states a proof-status limitation beyond the maximal-pants case. That is a precisely located later account. The unavailable joint chapter itself has not been directly verified, and the Earle--Sipe proposition above is not relabelled as a theorem from it.
